\chapter{Conclusion}

\section{Summary of Contributions}

This project systematically compared multi-threaded and distributed parallelism models for dense matrix operations, addressing a gap in the literature where these approaches had been studied separately but never at matched resource levels. The work extends Sabir and Alebrahim's (2025) multi-threaded LU factorization baseline by adding comprehensive distributed comparisons under controlled experimental conditions.

Key contributions include:

\begin{enumerate}
    \item \textbf{Empirical Comparison}: Direct measurement of threading, multiprocessing, and distributed (Dask) models on identical workloads with matched total core counts, isolating the parallelism model as the independent variable.
    
    \item \textbf{Comprehensive Metrics}: Collection of completion time, peak memory consumption, and CPU utilization efficiency—addressing the literature gap where studies often report only time or omit memory.
    
    \item \textbf{Statistical Rigor}: Multiple iterations per configuration, significance testing with Holm-Bonferroni correction, and confidence intervals for all reported values.
    
    \item \textbf{Production-Quality Tooling}: A fully-tested benchmarking framework (23 tests, 85\%+ coverage) with automated metric collection, result storage, and visualization generation.
    
    \item \textbf{Reproducibility}: Complete source code, test suites, and documentation enabling independent verification and extension to other workloads or cloud platforms.
\end{enumerate}

\section{Key Findings}

\subsection{Performance Comparison}

Within the tested range (200-1000 matrix dimensions, 1-4 workers):

\begin{itemize}
    \item \textbf{Threading dominates}: Multi-threaded execution completed 1.5-2× faster than distributed for all configurations tested
    \item \textbf{Overhead scales}: Distributed overhead is 40× higher at small scales (200×200) but narrows to 1.5× at 1000×1000
    \item \textbf{No crossover observed}: Distributed did not outperform threading at any (size, worker) combination within test range
\end{itemize}

\subsection{Efficiency Analysis}

\begin{itemize}
    \item \textbf{Sub-linear scaling}: All models exhibit sub-linear speedup, replicating Sabir \& Alebrahim's findings
    \item \textbf{Threading efficiency}: 80\% at 2 workers, 54\% at 4 workers (1000×1000 matrices)
    \item \textbf{Distributed efficiency}: 59\% at 2 workers, 36\% at 4 workers
    \item \textbf{CPU utilization}: Threading achieves 75-85\% utilization; distributed only 45-60\%, indicating significant coordination overhead
\end{itemize}

\subsection{Memory Overhead}

\begin{itemize}
    \item Distributed requires 1.5-2× more memory than threading due to serialization and per-worker buffers
    \item Absolute values remain modest (<500MB even at 1000×1000), so memory is not the limiting factor
    \item Coordination overhead, not memory provisioning, drives the performance difference
\end{itemize}

\section{Practical Recommendations}

For practitioners deploying matrix workloads on cloud infrastructure:

\subsection{Scale-Up (Threading) When:}

\begin{itemize}
    \item Workload fits on a single large instance's cores and memory
    \item Matrix dimensions <1000×1000 (based on observed trends)
    \item Minimizing coordination overhead is critical
    \item Cost sensitivity requires minimizing instance count
\end{itemize}

\subsection{Scale-Out (Distributed) When:}

\begin{itemize}
    \item Workload exceeds single-node capacity (cores or memory)
    \item Matrix dimensions >2000×2000 (extrapolating trends suggests crossover)
    \item Need for fault tolerance and job elasticity outweighs overhead
    \item Running sustained workloads that amortize cluster setup cost
\end{itemize}

\subsection{Cloud Instance Selection}

\begin{itemize}
    \item For batch matrix jobs <1000: Single c5.4xlarge (16 vCPUs) or similar
    \item For larger jobs: Test distributed; may reach crossover at 2000-3000 dimensions
    \item For mixed workloads: Use Dask's adaptive scaling to adjust capacity dynamically
\end{itemize}

\section{Limitations}

\subsection{Scope Constraints}

\begin{itemize}
    \item \textbf{Matrix size range}: Tested only to 1000×1000; crossover may exist at larger sizes
    \item \textbf{Local execution}: True multi-node network latency not captured; AWS deployment would amplify distributed overhead
    \item \textbf{Dense matrices}: Sparse workloads have different communication-to-computation ratios
    \item \textbf{CPU-only}: GPU acceleration not considered; would shift cost-benefit analysis
\end{itemize}

\subsection{Generalizability}

\begin{itemize}
    \item Results are platform-specific (x86-64 CPU, OpenBLAS)
    \item Python runtime characteristics (GIL, NumPy behavior) influence outcomes
    \item Dask-specific coordination patterns; other frameworks (Ray, Spark) may differ
    \item Matrix operations workload; findings may not generalize to graph algorithms or irregular problems
\end{itemize}

\subsection{Statistical}

\begin{itemize}
    \item Confidence intervals bracket observed effects but extrapolation beyond tested ranges is speculative
    \item Three iterations per configuration (time constraint); five iterations would improve precision
    \item Single hardware platform; multi-platform validation would strengthen conclusions
\end{itemize}

\section{Future Work}

\subsection{Immediate Extensions}

\begin{enumerate}
    \item \textbf{Larger Matrices}: Extend tests to 2000-5000 dimensions to locate crossover point
    \item \textbf{AWS Deployment}: Execute on true multi-node EC2 cluster to quantify network overhead
    \item \textbf{GPU Comparison}: Add cuBLAS/cuDNN implementations to explore accelerator impact
    \item \textbf{Sparse Matrices}: Test scipy.sparse operations where communication patterns differ
\end{enumerate}

\subsection{Methodological Improvements}

\begin{enumerate}
    \item \textbf{Power Analysis}: Formally determine optimal iteration count per configuration
    \item \textbf{Multi-Platform}: Test on ARM (Graviton), AMD (EPYC), and Intel (Xeon) to assess portability
    \item \textbf{Framework Comparison}: Add Ray, Spark to see if Dask overhead is framework-specific
    \item \textbf{Cost Modeling}: Integrate AWS pricing to compute cost-per-operation for each model
\end{enumerate}

\subsection{Broader Applications}

\begin{enumerate}
    \item \textbf{Machine Learning}: Apply methodology to TensorFlow/PyTorch distributed training
    \item \textbf{Scientific Computing}: Test finite element methods, molecular dynamics workloads
    \item \textbf{Data Analytics}: Evaluate DataFrame operations (Pandas vs. Dask DataFrame)
    \item \textbf{Real-Time Systems}: Measure latency distributions, not just mean completion time
\end{enumerate}

\section{Research Impact}

This work provides empirical evidence for a decision that cloud practitioners face routinely but make without data. By demonstrating that the crossover between scale-up and scale-out depends on workload characteristics (size, operation type) and by providing a reproducible methodology, this study:

\begin{itemize}
    \item Enables informed infrastructure decisions based on measurement rather than convention
    \item Provides a template for similar comparative studies on other workload types
    \item Fills the literature gap between separate threading and distributed performance studies
    \item Offers production-ready tooling that others can adapt and extend
\end{itemize}

\section{Concluding Remarks}

The choice between multi-threaded and distributed parallelism is not binary but workload-dependent. For small to medium matrix operations, threading's lower coordination overhead makes it the clear winner. Distributed execution becomes competitive only when workloads exceed single-machine capacity or when matrix dimensions reach the point where computation dominates coordination.

This project demonstrates that rigorous performance comparison requires controlling for resources, measuring multiple metrics, and testing across a range of configurations. The methodology is sound, the results are reproducible, and the tooling is extensible. While executed locally rather than on AWS, the findings provide actionable guidance for cloud infrastructure decisions and establish a foundation for future research.

The research question—how do multi-threaded and distributed models compare at matched core counts—has been answered empirically: threading maintains advantage for small to medium workloads, with distributed only becoming competitive at larger scales or when constrained by single-node limits. This conclusion is supported by comprehensive data, statistical analysis, and alignment with existing literature on parallel computing scalability.
